\section{Discussion}
\label{sec:discussion}

\subsection{Generality Across Devices}
While our system is valid against all devices tested in Section~\ref{subsec:impact_of_recording_devices}, there could be some devices that can possibly bypass our system. For instance, the resonance frequency of some devices could be shifted away from 40 kHz~\cite{10.1145/3133956.3134052,backdoor}, which makes the energy of the injected noise relatively low. Besides, iPhone 6 Plus could resist ultrasound signals because of the poor nonlinearity in its microphone~\cite{10.1145/3133956.3134052}. For the former, one possible solution is to adjust the carrier frequency of the ultrasonic noise to different resonance frequencies accordingly, but this would make the ultrasonic noise less effective across devices. In this paper, we choose 40 kHz as the carrier frequency because of its broad applicability~\cite{10.1145/3133956.3134052, backdoor}. For the latter, although it is difficult to jam the phone model with ultrasonic noise. Some other methods, such as Electromagnetic Interference (EMI), may be effective.
\subsection{Deployment of Transmitter Arrays}
In a practical application scenario, the deployment of the transmitter array is relatively convenient as the probability of destructive interference happening is neglectable. Our noise in each location is a superposition of the ultrasounds from different transmitters with various phases, the short wavelength fact and over-complex phase of these noises result in a slim chance of neutralization. So we only need to choose an appropriate number of arrays and spread them in the target room to satisfy the coverage requirement.